\section{Заключение}
\label{sec:Chapter6} \index{Chapter6}

В рамках работы был разработан механизм передачи диагностической информации из
гостевой системы в эмулятор Stratovirt. Решение включает виртуальное устройство
LogDev, драйвер для среды UEFI, адаптацию отладочной библиотеки EDK2 и драйвер
для операционной системы Windows.

В первой части работы были рассмотрены альтернативные подходы к передаче
диагностических сообщений: последовательный порт, debugcon,
virtio-console/virtio-serial и кольцевой буфер eBPF. Анализ показал, что эти
механизмы полезны в отдельных сценариях, но не совмещают все свойства,
необходимые для рассматриваемой задачи: раннее логирование UEFI, запись из
нескольких источников, единый лог гостя и эмулятора, а также возможность
минимизировать накладные расходы при записи из пользовательского пространства.

Исследовательская часть работы обосновала выбор общей памяти и MPSC-кольцевого
буфера без блокировок. Разделение операций резервирования, копирования и
фиксации записи позволяет нескольким писателям размещать сообщения в буфере и
не открывать читателю частично скопированные данные. MMIO-регистры при этом
используются не как основной канал передачи каждого байта, а как интерфейс
настройки и уведомления виртуального устройства.

Практическая реализация поддерживает два режима работы. До инициализации
оперативной памяти сообщения могут передаваться через синхронный регистр
\texttt{SYNC\_WR}. После выделения буфера в памяти гостевой системы основной
поток сообщений передаётся через разделяемый кольцевой буфер, а эмулятор
считывает данные после уведомления через \texttt{KICK} или \texttt{FLUSH}.
Драйвер Windows предоставляет интерфейсы для записи из пользовательского режима
и режима ядра.

Экспериментальная часть показала неоднозначный результат. По задержке доставки
сообщения до хоста стандартный \texttt{virtio-serial} оказался быстрее
реализованного прототипа: медианная задержка составила около 68--70 мс против
примерно 335 мс у LogDev с буфером без блокировок. В то же время вызов функции
записи в гостевой системе у LogDev оказался дешевле: в проведённых измерениях
варианты LogDev обрабатывали до 20 тыс. вызовов записи в секунду, тогда как
\texttt{virtio-serial} находился в диапазоне около 3.4--4.1 тыс. вызовов в
секунду. Сравнение двух вариантов LogDev показало преимущество lock-free
буфера над eBPF-подобной организацией по задержке доставки сообщения до хоста.

Дальнейшее развитие решения может включать оптимизацию пути доставки сообщения
на стороне хоста, повторные измерения при 16 потоках без единичных выбросов,
поддержку дополнительных гостевых операционных систем, расширение модели
безопасности для недоверенных писателей, добавление тайм-аутов для
незавершённых записей и автоматизированный набор производительных тестов для
регулярной проверки изменений в реализации.
